A fast method is presented for computing stability charts of linear autonomous time-delay systems: complete charts of simpler systems are generated in fractions of a second on a desktop machine -- fast enough for interactive, near-real-time parameter exploration. The number of unstable characteristic roots is obtained from the argument principle by tracking the phase of the characteristic function along the imaginary axis; the winding integral is reformulated as an ordinary differential equation and solved by an adaptive high-order integrator, with the integrand evaluated exactly by forward-mode automatic differentiation. The adaptive frequency march makes practically unbounded integration ranges affordable, and the result is self-validating: the distance of the computed winding number from the nearest integer is a dimensionless error measure that scales with the requested tolerance, so a single tolerance setting transfers between systems without problem-specific tuning. During the same sweep, the characteristic root closest to the integration line -- typically the rightmost root -- is estimated as a by-product. This estimate turns the otherwise piecewise-constant root-counting map into an interpolable field, so a single chart simultaneously displays the number of unstable roots and the spectral gap of the stable domain; it drives the multi-dimensional bisection method to trace the true stability boundary at high resolution, locates the most robust parameter point via the largest inscribed circle, and can optionally be sharpened toward machine precision by higher-order corrections built on the same automatic derivatives. The characteristic function is extracted algorithmically from a simulation-ready right-hand side, including singular mass matrices (delay differential--algebraic systems). Retarded, neutral, distributed-delay, transcendental, high-dimensional and fractional-order case studies demonstrate the identical workflow, and comparisons with adaptive quadrature and semi-discretization quantify the efficiency. The method is available in the open-source Julia package \pkg{InterpolatedNyquist.jl}.